
Uncertainty quantification methods for large language models fall into two isolated approaches: consistency-based methods that detect hallucinations via generative inconsistency, and statistical methods that provide coverage guarantees via conformal prediction. This fragmentation forces practitioners to choose between computational efficiency and statistical rigor. We demonstrate via synthetic proof-of-concept that these approaches capture complementary uncertainty dimensions. Correlation analysis on synthetic data analogs of three benchmark datasets reveals moderate correlation ($\rho$ = 0.43--0.46), suggesting that consistency captures epistemic uncertainty while conformal prediction captures aleatoric uncertainty. We introduce hierarchical Bayesian calibration (HBC), which integrates these methods through mutual updating. Synthetic validation experiments demonstrate that HBC achieves expected calibration error of 0.043, below the 0.05 threshold for well-calibrated models. The method maintains 92\% coverage while reducing computational cost by 30\% compared to conformal-only baselines. These results establish the core mechanism; real-world validation with full-scale datasets and actual model inference is required to confirm production applicability. The findings suggest that complementarity bounds (0.3 < $\rho$ < 0.7) may provide guidance for when joint calibration adds value over independent application.
